\documentclass[10pt,a4paper]{article}
\usepackage[utf8]{inputenc}
\usepackage[T1]{fontenc}
\usepackage[spanish]{babel}
\usepackage{geometry}
\geometry{left=1.5cm, right=1.5cm, top=1.5cm, bottom=1.5cm}
\usepackage{hyperref}
\usepackage{enumitem}
\usepackage{titlesec}

\titleformat{\section}{\large\bfseries\uppercase}{}{0em}{}[\titlerule]
\titlespacing{\section}{0pt}{1ex}{1ex}

\begin{document}
\pagestyle{empty}

\begin{center}
    {\huge \textbf{Jair Molina Arce}} \\
    \vspace{2mm}
    {\Large \textbf{Ingeniero Mecánico y de Sistemas Embebidos}} \\
    \vspace{2mm}
    Querétaro, Qro., México \\
    ingjairmolina@gmail.com | 5652646108 | jairmolina.dev | \href{https://github.com/smookymolina}{github.com/smookymolina} | \href{https://linkedin.com/in/jair-molina-arce-4909622b2}{linkedin.com/in/jair-molina-arce-4909622b2}
\end{center}

\section*{Perfil Profesional}
Ingeniero Mecánico con experiencia transversal en diseño mecánico, desarrollo de hardware y firmware para sistemas embebidos. He trabajado en caracterización de bancos de prueba, simulación térmica/estructural y automatización con microcontroladores (ESP32, STM32, RP2040). Destaco por mi capacidad para integrar la parte mecánica con la electrónica y la programación, creando soluciones completas y robustas desde el concepto hasta la validación física.

\section*{Educación}
\begin{itemize}[leftmargin=*, noitemsep]
    \item \textbf{Maestría en Tecnologías Avanzadas} (2024 -- Presente), Instituto Politécnico Nacional (IPN). Línea de investigación: control de sistemas dinámicos e instrumentación.
    \item \textbf{Ingeniería Mecánica} (2017 -- 2023), Instituto Politécnico Nacional (IPN). Enfoque en diseño de bancos de prueba, térmica, instrumentación y sistemas embebidos.
    \item \textbf{Bachillerato Técnico en Informática} (2014 -- 2017), Colegio de Bachilleres Plantel 1 ``El Rosario''.
\end{itemize}

\section*{Habilidades Técnicas}
\begin{itemize}[leftmargin=*, noitemsep]
    \item \textbf{Diseño y Simulación:} SolidWorks, ANSYS, KiCad, MATLAB/Simulink, diseño de bancos de prueba.
    \item \textbf{Firmware y Hardware:} C/C++, MicroPython, ESP32, STM32, RP2350, FSM, sensores, actuadores, diseño de PCB (mixed-signal), osciloscopio.
    \item \textbf{Software y Datos:} Python, LabVIEW, SQL, Docker, Git, Excel Avanzado.
\end{itemize}

\section*{Experiencia Relevante}
\begin{itemize}[leftmargin=*, noitemsep]
    \item \textbf{Docente Universitario} (Ago. 2025 -- Presente).
    \item \textbf{Ingeniería y Desarrollo Integrado - Banco de Pruebas de Propulsión} (2022 -- 2024). Diseño mecánico, FEA, instrumentación, DAQ, programación ESP32/STM32.
\end{itemize}

\section*{Proyecto Destacado}
\begin{itemize}[leftmargin=*, noitemsep]
    \item \textbf{Jaguar MX (2026) - Sistema de Extracción de Aire Caliente para Gabinete de Telecomunicaciones.} Desarrollo completo que incluyó diseño térmico, hardware, esquemático/PCB en KiCad, y firmware en C++ con máquinas de estado y validación en hardware.
\end{itemize}

\section*{Proyectos Seleccionados}
\begin{itemize}[leftmargin=*, noitemsep]
    \item Sistema de Telemetría Inalámbrica y Control Adaptativo PID.
    \item Pipeline de Automatización de Análisis de Datos y Reportes.
    \item Impresora 3D DIY - Diseño y Construcción desde Cero.
\end{itemize}

\section*{Freelance e Idiomas}
\begin{itemize}[leftmargin=*, noitemsep]
    \item \textbf{Freelance:} Plataforma Web Full-Stack para Monitoreo IoT (2023 -- 2024), Instructor de Excel Avanzado (2026).
    \item \textbf{Idiomas:} Español nativo | Inglés avanzado con lectura técnica fluida.
\end{itemize}

\end{document}
